\magnification=\magstep1
\hsize=16truecm
\vsize=24truecm
\parindent=0pt
\baselineskip=14pt
\raggedright
\tolerance=10000
\hbadness=10000

\font\tenrm="Liberation Serif" at 10pt
\font\tenbf="Liberation Serif/B" at 10pt
\font\tenit="Liberation Serif/I" at 10pt
\rm

\def\boxnum#1{\hbox{\vrule\vbox{\hrule\hbox{\strut\kern5pt{\bf#1}\kern5pt}\hrule}\vrule}}

% заголовок задачи
\def\problem#1#2{\vskip0pt plus 60pt\penalty-400\bigskip
  \line{\boxnum{#1}\quad #2\hfil}\nobreak\smallskip}
% заголовок шага решения
\def\st#1{\medskip\noindent{\bf #1}\enspace\ignorespaces}
% ответ в рамке
\def\ans#1{\medskip\noindent\hbox{\vrule\vbox{\hrule\hbox{\strut\kern5pt{\bf Ответ:}\enspace$\displaystyle #1$\kern5pt}\hrule}\vrule}\par}

%%%%%%%%%%%%%%%%%%%%%%%%%%%%%%%%%%%%%%%%%%%%%%%%%%%%%%%%%%%%%%%%%%%%%%%%
\centerline{\bf Уравнения 2-го порядка, допускающие понижение порядка}
\bigskip

{\bf Памятка: как решать.}\quad Общее решение уравнения 2-го порядка
$F(x,y,y',y'')=0$ содержит \underbar{две} произвольные постоянные $C_1$, $C_2$.
Если в уравнении нет части переменных, его порядок можно понизить подходящей
заменой: тогда вместо одного уравнения 2-го порядка мы решаем два уравнения
1-го порядка --- по одному на каждую постоянную. Сначала смотрим, чего в
уравнении \underbar{нет}.
\medskip

{\bf Тип 1. Уравнение $y''=f(x)$.}\quad Просто интегрируем два раза:
$$y'=\int f(x)\,dx+C_1,\qquad y=\int y'\,dx+C_2.$$
{\it Пример.} $y''=12x^2-\sin x$. Тогда $y'=4x^3+\cos x+C_1$, и
$y=x^4+\sin x+C_1x+C_2$.
\medskip

{\bf Тип 2. Нет $y$: $F(x,y',y'')=0$.}\quad Замена
$$y'=p(x),\qquad y''=p'(x).$$
Получаем уравнение 1-го порядка $F(x,p,p')=0$ относительно $p(x)$. Находим
$p=p(x,C_1)$, затем интегрируем ещё раз: $y=\int p(x,C_1)\,dx+C_2$.
(Если в уравнении нет и $y'$, а младшая производная --- $y^{(k)}$, то
заменяем $p=y^{(k)}$.)
\smallskip
{\it Пример.} $x\,y''+y'=4x$ ($x>0$). Замена $y'=p$, $y''=p'$:
$x\,p'+p=4x$. Левая часть --- это $(x\,p)'$, поэтому $(x\,p)'=4x$,
$x\,p=2x^2+C_1$, $p=2x+{C_1\over x}$. Интегрируем:
$y=x^2+C_1\ln x+C_2$.
\medskip

{\bf Тип 3. Нет $x$: $F(y,y',y'')=0$.}\quad Теперь за независимую
переменную берём $y$, а $y'$ считаем функцией от $y$:
$$y'=p(y),\qquad y''={dp\over dx}={dp\over dy}\cdot{dy\over dx}
=p\,{dp\over dy}.$$
Получаем уравнение 1-го порядка $F\bigl(y,p,p\,{dp\over dy}\bigr)=0$ относительно
$p(y)$. Находим $p=p(y,C_1)$, затем решаем $y'=p(y,C_1)$ --- это уравнение
с разделяющимися переменными: ${dy\over p(y,C_1)}=dx$, откуда
$\int{dy\over p(y,C_1)}=x+C_2$.
\smallskip
{\it Пример.} $y''+y=0$, $y(0)=0$, $y'(0)=1$. Замена $y'=p(y)$,
$y''=p\,{dp\over dy}$: $p\,{dp\over dy}=-y$, то есть $p\,dp=-y\,dy$ и
$p^2=-y^2+C_1$. Постоянную находим сразу: при $x=0$ имеем $y=0$,
$p=1$, значит $C_1=1$ и $y'^2=1-y^2$. Так как $y'(0)=1>0$, берём знак
плюс: $y'=\sqrt{1-y^2}$, ${dy\over\sqrt{1-y^2}}=dx$,
$\arcsin y=x+C_2$. Из $y(0)=0$ получаем $C_2=0$, то есть $y=\sin x$.
\medskip

{\bf Тип 4 (кратко). Однородное относительно $y,y',y''$.}\quad Если при замене
$y,y',y''\to ty,ty',ty''$ уравнение умножается на степень $t$, то подходит
замена
$$y'=y\,z(x),\qquad y''=y\,(z'+z^2).$$
Она сводит уравнение к уравнению 1-го порядка на $z(x)$; потом
$y=C_2\,e^{\int z\,dx}$ (решение $y=0$ нужно проверить отдельно).
\smallskip
{\it Пример.} $y\,y''-y'^2=y^2$. Подставляем:
$y^2(z'+z^2)-y^2z^2=y^2$, то есть $z'=1$, $z=x+C_1$. Тогда
${y'\over y}=x+C_1$ и $y=C_2\,e^{x^2/2+C_1x}$.
\medskip

{\bf Тип 5 (кратко). Левая часть --- полная производная.}\quad Если
$F(x,y,y',y'')=\bigl(\Phi(x,y,y')\bigr)'$, то уравнение равносильно
$\Phi(x,y,y')=C_1$ --- порядок понижен без всякой замены. Полезные заготовки:
$$(y\,y')'=y\,y''+y'^2,\qquad (x\,y')'=x\,y''+y',\qquad
(\ln y')'={y''\over y'}.$$
{\it Пример.} $y\,y''+y'^2=2$, то есть $(y\,y')'=2$. Тогда $y\,y'=2x+C_1$,
$\bigl({y^2\over 2}\bigr)'=2x+C_1$ и $y^2=2x^2+2C_1x+2C_2$; после
переобозначения постоянных $y^2=2x^2+C_1x+C_2$.
\medskip

{\bf Потерянные решения.}\quad При замене $y'=p(y)$ мы часто делим на $p$.
Решение $p=0$, то есть $y=\hbox{const}$, при этом теряется: его нужно
подставить в исходное уравнение отдельно (часто оно уже входит в общее
решение при специальном $C_1$, но не всегда). То же самое при любом
делении на выражение, которое может обратиться в нуль.
\smallskip
{\bf Задача Коши.}\quad Нужны два условия: $y(x_0)=y_0$, $y'(x_0)=y_0'$.
Постоянные лучше определять \underbar{сразу}, как только они появились: первую
($C_1$) --- из $y'(x_0)$ сразу после первого интегрирования (в типе 3 из пары
$y_0,y_0'$), вторую --- в конце. Это избавляет от громоздких выражений.
\smallskip
{\bf Знак при извлечении корня.}\quad Из $p^2=g(y)$ получаем
$p=\pm\sqrt{g(y)}$. Знак выбираем по условию $y'(x_0)=y_0'$: он совпадает со
знаком $y_0'$ (и сохраняется, пока $p$ не обращается в нуль). Если
$y_0'=0$, знак приходится выбирать из дополнительных соображений.
\smallskip
{\bf Про интегралы.}\quad $\int{dy\over\sqrt{1-y^2}}=\arcsin y$;\
$\int{dp\over1+p^2}=\arctan p$;\ $\int{dp\over1-p^2}={1\over2}\ln\left|{1+p\over1-p}\right|$;\
$\int\tan x\,dx=-\ln|\cos x|$.

\vfil\eject

%%%%%%%%%%%%%%%%%%%%%%%%%%%%%%%%%%%%%%%%%%%%%%%%%%%%%%%%%%%%%%%%%%%%%%%%
\centerline{\bf Классная работа}
\bigskip

%---------------------------------------------------------------- 1
\problem{1}{$y''=x+\cos x$}
\st{Тип.} Правая часть зависит только от $x$ --- интегрируем два раза.
\st{Шаг 1.} Первое интегрирование:
$$y'=\int(x+\cos x)\,dx={x^2\over2}+\sin x+C_1.$$
\st{Шаг 2.} Второе интегрирование:
$$y=\int\left({x^2\over2}+\sin x+C_1\right)dx={x^3\over6}-\cos x+C_1x+C_2.$$
\ans{y={x^3\over6}-\cos x+C_1x+C_2}

%---------------------------------------------------------------- 2
\problem{2}{$y''={1\over\cos^2x},\quad y(0)=0,\ y'(0)=1$}
\st{Тип.} $y''=f(x)$, задача Коши.
\st{Шаг 1.} Интегрируем: $y'=\tan x+C_1$. Условие $y'(0)=1$ даёт
$0+C_1=1$, то есть $C_1=1$ и $y'=\tan x+1$.
\st{Шаг 2.} Интегрируем ещё раз (используем $\int\tan x\,dx=-\ln|\cos x|$):
$$y=-\ln|\cos x|+x+C_2.$$
Условие $y(0)=0$: $-\ln1+0+C_2=0$, то есть $C_2=0$.
\smallskip
Решение определено при $|x|<{\pi\over2}$, там $\cos x>0$.
\ans{y=x-\ln\cos x}

%---------------------------------------------------------------- 3
\problem{3}{$x\,y''=y'$}
\st{Тип.} В уравнении нет $y$. Замена $y'=p(x)$, $y''=p'$.
\st{Шаг 1.} Получаем $x\,p'=p$ (считаем $x>0$). Разделяем переменные:
$p=0$ или
$${dp\over p}={dx\over x}\quad\Rightarrow\quad\ln|p|=\ln x+c
\quad\Rightarrow\quad p=C_1x.$$
Решение $p=0$ входит сюда при $C_1=0$.
\st{Шаг 2.} Возвращаемся к $y$: $y'=C_1x$, откуда
$y={C_1x^2\over2}+C_2$. Постоянная $C_1/2$ --- тоже произвольная, обозначим
её снова $C_1$.
\ans{y=C_1x^2+C_2}

%---------------------------------------------------------------- 4
\problem{4}{$y''-y'=x$}
\st{Тип.} Нет $y$. Замена $y'=p(x)$, $y''=p'$:
$$p'-p=x.$$
Это линейное уравнение 1-го порядка ($P=-1$, $Q=x$).
\st{Шаг 1.} Однородное $p'-p=0$ даёт $p=C\,e^{x}$. Варьируем постоянную:
$p=C(x)\,e^{x}$, $p'=C'e^x+C\,e^x$. Подставляем в $p'-p=x$:
$$C'e^{x}+C\,e^x-C\,e^x=x\quad\Rightarrow\quad C'=x\,e^{-x}.$$
По частям: $\int x\,e^{-x}dx=-x\,e^{-x}-e^{-x}$, поэтому
$C(x)=-(x+1)\,e^{-x}+C_1$ и
$$p=C(x)\,e^{x}=-x-1+C_1e^{x}.$$
Проверка: $p'-p=C_1e^x-1+x+1-C_1e^x=x$.
\st{Шаг 2.} Возвращаемся к $y$: $y'=-x-1+C_1e^x$, откуда
$$y=-{x^2\over2}-x+C_1e^{x}+C_2.$$
\ans{y=-{x^2\over2}-x+C_1e^{x}+C_2}

%---------------------------------------------------------------- 5
\problem{5}{$(1+x^2)\,y''=2x\,y',\quad y(0)=1,\ y'(0)=3$}
\st{Тип.} Нет $y$, задача Коши. Замена $y'=p(x)$, $y''=p'$:
$(1+x^2)\,p'=2x\,p$.
\st{Шаг 1.} Разделяем переменные ($p=0$ --- решение, оно войдёт при
$C_1=0$):
$${dp\over p}={2x\,dx\over1+x^2}\quad\Rightarrow\quad
\ln|p|=\ln(1+x^2)+c\quad\Rightarrow\quad p=C_1(1+x^2).$$
Сразу определяем $C_1$ из $y'(0)=p(0)=3$: $C_1\cdot1=3$, то есть $C_1=3$ и
$$y'=3(1+x^2)=3+3x^2.$$
\st{Шаг 2.} Интегрируем: $y=3x+x^3+C_2$. Из $y(0)=1$ получаем $C_2=1$.
\smallskip
(Общее решение: $y=C_1\left(x+{x^3\over3}\right)+C_2$.)
\ans{y=x^3+3x+1}

%---------------------------------------------------------------- 6
\problem{6}{$y\,y''=y'^2$}
\st{Тип.} Нет $x$. Замена $y'=p(y)$, $y''=p\,{dp\over dy}$:
$$y\,p\,{dp\over dy}=p^2.$$
\st{Шаг 1.} Случай $p=0$: $y'=0$, $y=\hbox{const}$ --- это решения (проверяем
подстановкой: $y\cdot0=0$). Пусть $p\ne0$; сокращаем на $p$ и разделяем
переменные:
$${dp\over p}={dy\over y}\quad\Rightarrow\quad\ln|p|=\ln|y|+c
\quad\Rightarrow\quad p=C_1y.$$
(При $C_1=0$ получаем снова $p=0$.)
\st{Шаг 2.} Возвращаемся к $y$: $y'=C_1y$. Это уравнение с разделяющимися
переменными:
$${dy\over y}=C_1\,dx\quad\Rightarrow\quad\ln|y|=C_1x+c
\quad\Rightarrow\quad y=C_2\,e^{C_1x}.$$
При $C_1=0$ имеем $y=C_2=\hbox{const}$, так что потерянные решения
входят в ответ.
\smallskip
(Это же уравнение однородное: $y'=y\,z$ даёт $z'=0$, то есть $z=C_1$ --- тот же
ответ.)
\ans{y=C_2\,e^{C_1x}}

%---------------------------------------------------------------- 7
\problem{7}{$y''=2y\,y',\quad y(0)=0,\ y'(0)=1$}
\st{Тип.} Нет $x$. Замена $y'=p(y)$, $y''=p\,{dp\over dy}$:
$$p\,{dp\over dy}=2y\,p.$$
\st{Шаг 1.} Решение $p=0$ (то есть $y=\hbox{const}$) нам не подходит,
так как $y'(0)=1\ne0$. При $p\ne0$ сокращаем на $p$:
$${dp\over dy}=2y\quad\Rightarrow\quad p=y^2+C_1.$$
Определяем $C_1$: при $x=0$ имеем $y=0$, $p=1$, значит $C_1=1$ и
$$y'=y^2+1.$$
\st{Шаг 2.} Разделяем переменные:
$${dy\over1+y^2}=dx\quad\Rightarrow\quad\arctan y=x+C_2.$$
Из $y(0)=0$ получаем $C_2=0$, то есть $y=\tan x$ (при $|x|<{\pi\over2}$).
\smallskip
(Заметим, что уравнение --- это $y''=(y^2)'$, то есть тип 5: сразу
$y'=y^2+C_1$.)
\ans{y=\tan x}

%---------------------------------------------------------------- 8
\problem{8}{$y''=2y^3,\quad y(0)=1,\ y'(0)=1$}
\st{Тип.} Нет $x$. Замена $y'=p(y)$, $y''=p\,{dp\over dy}$:
$p\,{dp\over dy}=2y^3$.
\st{Шаг 1.} Разделяем переменные в виде $p\,dp=2y^3dy$:
$${p^2\over2}={y^4\over2}+c\quad\Rightarrow\quad p^2=y^4+C_1.$$
Определяем $C_1$: при $x=0$ имеем $y=1$, $p=1$, значит
$1=1+C_1$, $C_1=0$. Тогда $p^2=y^4$, то есть $p=\pm y^2$. Так как
$y'(0)=1>0$, берём знак плюс:
$$y'=y^2.$$
\st{Шаг 2.} Разделяем переменные:
$${dy\over y^2}=dx\quad\Rightarrow\quad-{1\over y}=x+C_2.$$
Из $y(0)=1$: $-1=C_2$. Значит $-{1\over y}=x-1$, то есть $y={1\over1-x}$
(при $x<1$).
\smallskip
Проверка: $y'={1\over(1-x)^2}=y^2$, $y''={2\over(1-x)^3}=2y^3$ --- верно.
\ans{y={1\over1-x}}

%---------------------------------------------------------------- 9
\problem{9}{$y''=1-y'^2$}
\st{Тип.} Нет ни $x$, ни $y$ --- годятся обе замены; берём $y'=p(x)$,
$y''=p'$:
$$p'=1-p^2.$$
\st{Шаг 1.} Потерянные решения: $p=1$ и $p=-1$ (правая часть обращается в
нуль). Они дают $y'=\pm1$, то есть $y=\pm x+C$; проверка: $y''=0=1-1$ ---
верно. При $p\ne\pm1$ разделяем переменные:
$${dp\over1-p^2}=dx\quad\Rightarrow\quad{1\over2}\ln\left|{1+p\over1-p}\right|=x+C_1.$$
Выражая $p$: ${1+p\over1-p}=\pm e^{2(x+C_1)}$. При знаке плюс получаем
($|p|<1$)
$$p=\tanh(x+C_1)={e^{2(x+C_1)}-1\over e^{2(x+C_1)}+1},$$
при знаке минус --- $p=\coth(x+C_1)$ ($|p|>1$).
\st{Шаг 2.} Интегрируем $y'=p$ (используем $\int\tanh t\,dt=\ln\cosh t$,
$\int\coth t\,dt=\ln|\sinh t|$):
$$y=\ln\cosh(x+C_1)+C_2\quad\hbox{или}\quad y=\ln|\sinh(x+C_1)|+C_2.$$
Общий ответ также включает $y=\pm x+C$.
\ans{y=\ln\cosh(x+C_1)+C_2,\quad y=\ln|\sinh(x+C_1)|+C_2,\quad
y=\pm x+C}

%---------------------------------------------------------------- 10
\problem{10}{$y^3\,y''=1,\quad y(0)=1,\ y'(0)=1$}
\st{Тип.} Нет $x$. Замена $y'=p(y)$, $y''=p\,{dp\over dy}$:
$$y^3\,p\,{dp\over dy}=1\quad\Rightarrow\quad p\,dp={dy\over y^3}.$$
\st{Шаг 1.} Интегрируем:
$${p^2\over2}=-{1\over2y^2}+c\quad\Rightarrow\quad p^2=C_1-{1\over y^2}.$$
Определяем $C_1$: при $x=0$ имеем $y=1$, $p=1$, значит $1=C_1-1$,
$C_1=2$. Тогда
$$y'^2={2y^2-1\over y^2}.$$
Так как $y'(0)=1>0$ и $y>0$, берём знак плюс:
$y'={\sqrt{2y^2-1}\over y}$.
\st{Шаг 2.} Разделяем переменные:
$${y\,dy\over\sqrt{2y^2-1}}=dx.$$
Для $u=2y^2-1$ имеем $du=4y\,dy$, поэтому
$\int{y\,dy\over\sqrt{2y^2-1}}={1\over4}\int{du\over\sqrt u}={\sqrt{2y^2-1}\over2}$.
Значит
$${\sqrt{2y^2-1}\over2}=x+C_2.$$
Из $y(0)=1$: $C_2={1\over2}$. Тогда $\sqrt{2y^2-1}=2x+1$ (при $x>-{1\over2}$);
возводим в квадрат: $2y^2-1=4x^2+4x+1$, откуда $y^2=2x^2+2x+1$. Так как
$y(0)=1>0$, берём положительный корень.
\smallskip
Проверка: $y^2=Q:=2x^2+2x+1$, $y\,y'={Q'\over2}=2x+1$, далее
$y'^2+y\,y''=2$, откуда $y\,y''=2-{(2x+1)^2\over Q}={1\over Q}$ и
$y^3y''=Q\cdot{1\over Q}=1$.
\ans{y=\sqrt{2x^2+2x+1}}

\vfil\eject

%%%%%%%%%%%%%%%%%%%%%%%%%%%%%%%%%%%%%%%%%%%%%%%%%%%%%%%%%%%%%%%%%%%%%%%%
\centerline{\bf Домашняя работа}
\bigskip

\line{\boxnum{1}\quad $y''=x\,e^{x}$\hfil}
\bigskip

\line{\boxnum{2}\quad $y''={1\over x},\quad y(1)=0,\ y'(1)=1$\hfil}
\bigskip

\line{\boxnum{3}\quad $y''\tan x=y'$\hfil}
\bigskip

\line{\boxnum{4}\quad $y''+2y'=4x$\hfil}
\bigskip

\line{\boxnum{5}\quad $x\,y''=y'+x^2,\quad y(1)=1,\ y'(1)=3$\hfil}
\bigskip

\line{\boxnum{6}\quad $y\,y''+y'^2=6x$\hfil}
\bigskip

\line{\boxnum{7}\quad $y''=e^{2y},\quad y(0)=0,\ y'(0)=1$\hfil}
\bigskip

\line{\boxnum{8}\quad $y\,y''=2y'^2,\quad y(0)=1,\ y'(0)=-2$\hfil}
\bigskip

\line{\boxnum{9}\quad $y''=1+y'^2,\quad y(0)=0,\ y'(0)=0$\hfil}
\bigskip

\line{\boxnum{10}\quad $y\,y''-y'^2=x\,y^2,\quad y(0)=1,\ y'(0)=2$\hfil}

\vfil\eject

%%%%%%%%%%%%%%%%%%%%%%%%%%%%%%%%%%%%%%%%%%%%%%%%%%%%%%%%%%%%%%%%%%%%%%%%
\centerline{\bf Домашняя работа: решения}
\bigskip

%---------------------------------------------------------------- H1
\problem{1}{$y''=x\,e^{x}$}
\st{Тип.} $y''=f(x)$.
\st{Шаг 1.} По частям ($p=x$, $dq=e^xdx$):
$\int x\,e^xdx=x\,e^x-e^x$. Значит $y'=(x-1)\,e^x+C_1$.
\st{Шаг 2.} Ещё раз по частям:
$\int(x-1)\,e^xdx=(x-1)\,e^x-\int e^xdx=(x-2)\,e^x$. Значит
$$y=(x-2)\,e^x+C_1x+C_2.$$
\ans{y=(x-2)\,e^{x}+C_1x+C_2}

%---------------------------------------------------------------- H2
\problem{2}{$y''={1\over x},\quad y(1)=0,\ y'(1)=1$}
\st{Тип.} $y''=f(x)$, задача Коши ($x>0$).
\st{Шаг 1.} $y'=\ln x+C_1$. Из $y'(1)=1$: $0+C_1=1$, $C_1=1$,
$y'=\ln x+1$.
\st{Шаг 2.} $\int\ln x\,dx=x\ln x-x$, поэтому
$y=x\ln x-x+x+C_2=x\ln x+C_2$. Из $y(1)=0$: $C_2=0$.
\ans{y=x\ln x}

%---------------------------------------------------------------- H3
\problem{3}{$y''\tan x=y'$}
\st{Тип.} Нет $y$. Замена $y'=p(x)$, $y''=p'$: $p'\tan x=p$.
\st{Шаг 1.} Решение $p=0$ войдёт в общее при $C_1=0$. При $p\ne0$:
$${dp\over p}=\cot x\,dx\quad\Rightarrow\quad\ln|p|=\ln|\sin x|+c
\quad\Rightarrow\quad p=C_1\sin x.$$
\st{Шаг 2.} $y'=C_1\sin x$, откуда $y=-C_1\cos x+C_2$. Обозначим
$-C_1$ снова через $C_1$.
\ans{y=C_1\cos x+C_2}

%---------------------------------------------------------------- H4
\problem{4}{$y''+2y'=4x$}
\st{Тип.} Нет $y$. Замена $y'=p(x)$, $y''=p'$: $p'+2p=4x$ --- линейное
уравнение.
\st{Шаг 1.} Однородное $p'+2p=0$ даёт $p=C\,e^{-2x}$. Варьируем
постоянную: $p=C(x)\,e^{-2x}$, $p'=C'e^{-2x}-2C\,e^{-2x}$. Подставляем:
$$C'e^{-2x}-2C\,e^{-2x}+2C\,e^{-2x}=4x\quad\Rightarrow\quad C'=4x\,e^{2x}.$$
По частям: $\int4x\,e^{2x}dx=2x\,e^{2x}-\int2e^{2x}dx=(2x-1)\,e^{2x}$. Значит
$C(x)=(2x-1)\,e^{2x}+C_1$ и
$$p=2x-1+C_1e^{-2x}.$$
Проверка: $p'+2p=2-2C_1e^{-2x}+4x-2+2C_1e^{-2x}=4x$.
\st{Шаг 2.} Интегрируем $y'=p$:
$y=x^2-x-{C_1\over2}e^{-2x}+C_2$. Обозначим $-{C_1\over2}$ снова через $C_1$.
\ans{y=x^2-x+C_1e^{-2x}+C_2}

%---------------------------------------------------------------- H5
\problem{5}{$x\,y''=y'+x^2,\quad y(1)=1,\ y'(1)=3$}
\st{Тип.} Нет $y$, задача Коши ($x>0$). Замена $y'=p(x)$, $y''=p'$:
$x\,p'=p+x^2$. Делим на $x$:
$$p'-{p\over x}=x.$$
Это линейное уравнение.
\st{Шаг 1.} Однородное $p'-{p\over x}=0$ даёт $p=C\,x$. Варьируем
постоянную: $p=C(x)\,x$, $p'=C'x+C$. Подставляем:
$$C'x+C-C=x\quad\Rightarrow\quad C'=1\quad\Rightarrow\quad C(x)=x+C_1.$$
Значит, $p=x^2+C_1x$. Из $y'(1)=p(1)=3$: $1+C_1=3$, $C_1=2$, то есть
$y'=x^2+2x$.
\st{Шаг 2.} Интегрируем: $y={x^3\over3}+x^2+C_2$. Из $y(1)=1$:
${1\over3}+1+C_2=1$, $C_2=-{1\over3}$.
\ans{y={x^3\over3}+x^2-{1\over3}}

%---------------------------------------------------------------- H6
\problem{6}{$y\,y''+y'^2=6x$}
\st{Тип.} Слева стоит полная производная: $(y\,y')'=y\,y''+y'^2$. Поэтому
уравнение имеет вид $(y\,y')'=6x$.
\st{Шаг 1.} Интегрируем: $y\,y'=3x^2+C_1$.
\st{Шаг 2.} Замечаем, что $y\,y'=\bigl({y^2\over2}\bigr)'$. Интегрируем ещё раз:
$${y^2\over2}=x^3+C_1x+C_2\quad\Rightarrow\quad y^2=2x^3+2C_1x+2C_2.$$
После переобозначения постоянных получаем ответ. Он остаётся в неявной
форме ($y=\pm\sqrt{\ldots}$, знак выбирается по условию).
\smallskip
Проверка: $y\,y'=3x^2+C_1$, $y'^2+y\,y''=6x$ --- верно.
\ans{y^2=2x^3+C_1x+C_2}

%---------------------------------------------------------------- H7
\problem{7}{$y''=e^{2y},\quad y(0)=0,\ y'(0)=1$}
\st{Тип.} Нет $x$, задача Коши. Замена $y'=p(y)$, $y''=p\,{dp\over dy}$:
$p\,{dp\over dy}=e^{2y}$.
\st{Шаг 1.} Разделяем переменные: $p\,dp=e^{2y}dy$, откуда
$${p^2\over2}={e^{2y}\over2}+c\quad\Rightarrow\quad p^2=e^{2y}+C_1.$$
Определяем $C_1$: при $x=0$ имеем $y=0$, $p=1$, значит $1=1+C_1$, $C_1=0$.
Тогда $p^2=e^{2y}$, $p=\pm e^{y}$; так как $y'(0)=1>0$, берём знак плюс:
$y'=e^{y}$.
\st{Шаг 2.} Разделяем переменные:
$$e^{-y}dy=dx\quad\Rightarrow\quad-e^{-y}=x+C_2.$$
Из $y(0)=0$: $-1=C_2$. Значит, $e^{-y}=1-x$, и
$y=-\ln(1-x)$ (при $x<1$).
\smallskip
Проверка: $y'={1\over1-x}=e^{y}$, $y''={1\over(1-x)^2}=e^{2y}$ --- верно.
\ans{y=-\ln(1-x)}

%---------------------------------------------------------------- H8
\problem{8}{$y\,y''=2y'^2,\quad y(0)=1,\ y'(0)=-2$}
\st{Тип.} Нет $x$ (и уравнение однородное), задача Коши. Замена $y'=p(y)$,
$y''=p\,{dp\over dy}$:
$$y\,p\,{dp\over dy}=2p^2.$$
\st{Шаг 1.} Решение $p=0$ не подходит, так как $y'(0)=-2\ne0$. Сокращаем
на $p$ и разделяем переменные:
$${dp\over p}=2\,{dy\over y}\quad\Rightarrow\quad\ln|p|=2\ln|y|+c
\quad\Rightarrow\quad p=C_1y^2.$$
Определяем $C_1$: при $x=0$ имеем $y=1$, $p=-2$, значит $C_1=-2$ и
$$y'=-2y^2.$$
\st{Шаг 2.} Разделяем переменные:
$${dy\over y^2}=-2\,dx\quad\Rightarrow\quad-{1\over y}=-2x+C_2.$$
Из $y(0)=1$: $C_2=-1$. Значит, ${1\over y}=2x+1$ и $y={1\over1+2x}$
(при $x>-{1\over2}$).
\smallskip
(Общее решение: $y={1\over C_1x+C_2}$ и $y=\hbox{const}$; константа с
$C_1=0$ входит в формулу.)
\ans{y={1\over1+2x}}

%---------------------------------------------------------------- H9
\problem{9}{$y''=1+y'^2,\quad y(0)=0,\ y'(0)=0$}
\st{Тип.} Нет $y$ (и нет $x$), задача Коши. Замена $y'=p(x)$, $y''=p'$:
$p'=1+p^2$. Правая часть не обращается в нуль --- решения не теряются.
\st{Шаг 1.} Разделяем переменные:
$${dp\over1+p^2}=dx\quad\Rightarrow\quad\arctan p=x+C_1.$$
Из $p(0)=y'(0)=0$: $C_1=0$, значит $p=\tan x$, то есть $y'=\tan x$.
\st{Шаг 2.} $y=-\ln|\cos x|+C_2$. Из $y(0)=0$: $C_2=0$. Решение определено при
$|x|<{\pi\over2}$.
\smallskip
(Общее решение: $y=-\ln|\cos(x+C_1)|+C_2$.)
\ans{y=-\ln\cos x}

%---------------------------------------------------------------- H10
\problem{10}{$y\,y''-y'^2=x\,y^2,\quad y(0)=1,\ y'(0)=2$}
\st{Тип.} Уравнение однородное относительно $y,y',y''$ (при замене
$y\to ty$ обе части умножаются на $t^2$). Замена
$$y'=y\,z(x),\qquad y''=y'z+y\,z'=y\,(z^2+z').$$
\st{Шаг 1.} Подставляем в уравнение:
$$y\cdot y\,(z^2+z')-y^2z^2=x\,y^2\quad\Rightarrow\quad y^2z'=x\,y^2
\quad\Rightarrow\quad z'=x.$$
(Делим на $y^2$; решение $y=0$ войдёт в ответ при $C_2=0$.)
Интегрируем: $z={x^2\over2}+C_1$. Определяем $C_1$: $z(0)={y'(0)\over y(0)}=2$,
значит $C_1=2$.
\st{Шаг 2.} Возвращаемся к $y$: ${y'\over y}={x^2\over2}+2$,
$$\ln|y|={x^3\over6}+2x+c\quad\Rightarrow\quad y=C_2\,e^{x^3/6+2x}.$$
Из $y(0)=1$: $C_2=1$.
\smallskip
(Общее решение: $y=C_2\,e^{x^3/6+C_1x}$.)
\ans{y=e^{x^3/6+2x}}

\bye
